\chapter{Технологическая часть}

\section{Средства разработки}

В качестве языка программирования был выбран MicroPython \cite{micropython}, так как он поддерживается микроконтроллером, на котором будут выполняться замеры, в его стандартной библиотеке присутствуют функция замера процессорного времени, которая требуется в условии.

Результаты замеров сохраняются в .csv файл, для построения графиков использовался exel.

Для замера времени использовалась функция ticks\_ms() из стандартного модуля utime \cite{micropython-utime}.

\section{Реализация алгоритмов}

В листинге \ref{lst:std} представлена реализация стандартного алгоритма умножения матриц.

\begin{lstlisting}[label=lst:std,caption={Стандартный алгоритм умножения матриц}]
def Multiply(A, B):
    rows_A = len(A)
    cols_A = len(A[0])
    rows_B = len(B)
    cols_B = len(B[0])

    result = [[0 for _ in range(cols_B)] for _ in range(rows_A)]

    for i in range(rows_A):
        for j in range(cols_B):
            for k in range(cols_A):
                result[i][j] += A[i][k] * B[k][j]

    return result
\end{lstlisting}

\clearpage

В листинге \ref{lst:vin} представлена реализация алгоритма умножения матриц Винограда.

\begin{lstlisting}[label=lst:vin,caption={Алгоритм умножения матриц Винограда}]
def VinogradMultiply(A, B):
    rows_A = len(A)
    cols_A = len(A[0])
    rows_B = len(B)
    cols_B = len(B[0])

    result = [[0 for _ in range(cols_B)] for _ in range(rows_A)]

    mulH = [0] * rows_A
    for i in range(rows_A):
        for j in range(0, cols_A // 2):
            mulH[i] = mulH[i] + A[i][2 * j] * A[i][2 * j + 1]

    mulV = [0] * cols_B
    for i in range(cols_B):
        for j in range(0, rows_B // 2):
            mulV[i] = mulV[i] + B[2 * j][i] * B[2 * j + 1][i]

    for i in range(rows_A):
        for j in range(cols_B):
            result[i][j] = -mulH[i] - mulV[j]
            for k in range(0, cols_A // 2):
                result[i][j] = result[i][j] + (A[i][2 * k] + B[2 * k + 1][j]) * (A[i][2 * k + 1] + B[2 * k][j])

    if cols_A % 2 == 1:
        for i in range(rows_A):
            for j in range(cols_B):
                result[i][j] = result[i][j] + A[i][cols_A - 1] * B[cols_A - 1][j]

    return result
\end{lstlisting}

\clearpage

В листинге \ref{lst:opt-vin} представлена реализация оптимизированного алгоритма умножения матриц Винограда.
\begin{itemize}
    \item двоичный сдвиг вместо умножения на 2
    \item введение декремента при вычислении вспомогательных массивов
    \item вынос начальной итерации из каждого внешнего цикл
\end{itemize}

\begin{lstlisting}[label=lst:opt-vin,caption={Алгоритм умножения матриц Винограда}]
def OptimVinogradMultiply(A, B):
    rows_A = len(A)
    cols_A = len(A[0])
    rows_B = len(B)
    cols_B = len(B[0])

    result = [[0 for _ in range(cols_B)] for _ in range(rows_A)]
    mulH = [0] * rows_A
    for i in range(rows_A):
        mulH[i] -= A[i][0] * A[i][1]
        for j in range(1, cols_A // 2):
            mulH[i] -= A[i][j << 1] * A[i][(j << 1 )+ 1]
    mulV = [0] * cols_B
    for i in range(cols_B):
        mulV[i] -= B[0][i] * B[1][i]
        for j in range(1, rows_B // 2):
            mulV[i] -= B[j << 1][i] * B[(j << 1) + 1][i]
    for i in range(rows_A):
        for j in range(cols_B):
            result[i][j] = mulH[i] + mulV[j]
            result[i][j] = result[i][j] + (A[i][0] + B[1][j]) * (A[i][ 1] + B[0][j])
            for k in range(1, cols_A // 2):
                result[i][j] = result[i][j] + (A[i][k << 1] + B[(k << 1) + 1][j]) * (A[i][(k << 1) + 1] + B[k << 1][j])
    if cols_A % 2 == 1:
        for i in range(rows_A):
            for j in range(cols_B):
                result[i][j] = result[i][j] + A[i][cols_A - 1] * B[cols_A - 1][j]
    return result
\end{lstlisting}

\clearpage

\section{Функциональные тесты}

В табличке \ref{tbl:func} приведены функциональные тесты для алгоритмов.

\begin{table}[H]
\begin{tabularx}{\textwidth} { 
  | >{\centering\arraybackslash}X 
  | >{\centering\arraybackslash}X 
  | >{\centering\arraybackslash}X 
  | >{\centering\arraybackslash}X | }
        \hline
        \textbf{№} & \textbf{Входные данные} & \textbf{Результат} & \textbf{Ожидаемое} \\ 
        \hline
        1 & 
        $A = \begin{pmatrix} 
        1 & 2 & 3 \\ 
        4 & 5 & 6 \\ 
        7 & 8 & 9 
        \end{pmatrix}$ \newline 
        $B = \begin{pmatrix} 
        9 & 8 & 7 \\ 
        6 & 5 & 4 \\ 
        3 & 2 & 1 
        \end{pmatrix}$ & 
        $C = \begin{pmatrix} 
        30 & 24 & 18 \\ 
        84 & 69 & 54 \\ 
        138 & 114 & 90 
        \end{pmatrix}$ & 
        Результат совпадает с рассчитанной матрицей C. \\ 
        \hline
        
        2 & 
        $A = \begin{pmatrix} 
        1 & 2 & 3 \\ 
        4 & 5 & 6 
        \end{pmatrix}$ \newline 
        $B = \begin{pmatrix} 
        7 & 8 \\ 
        9 & 10 \\ 
        11 & 12 
        \end{pmatrix}$ & 
        $C = \begin{pmatrix} 
        58 & 64 \\ 
        139 & 154 
        \end{pmatrix}$ & 
        Матрица C совпадает с ожидаемой. \\ 
        \hline
        
        3 & 
        $A = \begin{pmatrix} 
        1 & 2 & 3 & 4 \\ 
        0 & 1 & 2 & 3 \\ 
        0 & 0 & 1 & 2 \\ 
        0 & 0 & 0 & 1 
        \end{pmatrix}$ \newline 
        $I = \begin{pmatrix} 
        1 & 0 & 0 & 0 \\ 
        0 & 1 & 0 & 0 \\ 
        0 & 0 & 1 & 0 \\ 
        0 & 0 & 0 & 1 
        \end{pmatrix}$ & 
        $A = \begin{pmatrix} 
        1 & 2 & 3 & 4 \\ 
        0 & 1 & 2 & 3 \\ 
        0 & 0 & 1 & 2 \\ 
        0 & 0 & 0 & 1 
        \end{pmatrix}$ & 
        Результат совпадает с исходной матрицей A. \\ 
        \hline
\end{tabularx}
\caption{Функциональные тесты для алгоритма умножения матриц с конкретными данными}
\label{tbl:func}
\end{table}


Все тесты пройдены успешно

\section{Вывод}

В ходе работы были разработаны стандартный алгоритм и алгоритм Винограда умножения матриц и оптимизированный по варианту алгоритм Винограда, а также проведено их тестирование.

\clearpage
